\clearpage\thispagestyle{empty}
\pdfbookmark[0]{Quellen zu Pallas (24. Februar 1810)}{pallas1810feb24-sources}
\begin{center}\textbf{Quellen zu Pallas (24. Februar 1810)}\end{center}
\noindent Diese KI-Übertragung ins heutige Deutsch folgt dem Bericht in \href{https://gdz.sub.uni-goettingen.de/id/PPN236007467}{Gauß’ \textit{Werken}, Bd.~VI} (Göttingen, 1874), S.~317--318. Er erschien in den \textit{Göttingischen gelehrten Anzeigen}, Stück 32, 24. Februar 1810, S.~\href{https://api.digitale-sammlungen.de/iiif/image/v2/bsb10538677_00341/full/1600,/0/default.jpg}{305}, \href{https://api.digitale-sammlungen.de/iiif/image/v2/bsb10538677_00342/full/1600,/0/default.jpg}{306}, \href{https://api.digitale-sammlungen.de/iiif/image/v2/bsb10538677_00343/full/1600,/0/default.jpg}{307}, \href{https://api.digitale-sammlungen.de/iiif/image/v2/bsb10538677_00344/full/1600,/0/default.jpg}{308}. Das digitalisierte Exemplar stammt aus der Bayerischen Staatsbibliothek.
\par\smallskip\noindent Die Werkausgabe gibt die nördliche Breite bei der Opposition von 1808 mit 37\,33\,54 an. In beiden Zeitschriftenberichten steht 37\,43\,54 (siehe \href{https://api.digitale-sammlungen.de/iiif/image/v2/bsb10538677_00342/full/1600,/0/default.jpg}{S.~306} und \href{https://api.digitale-sammlungen.de/iiif/image/v2/bsb10538611_00285/full/1600,/0/default.jpg}{\textit{Monatliche Correspondenz}, S.~277}). In der Werkausgabe und in den \textit{Anzeigen} endet die Deklination vom 29. September auf 57.3, in der \href{https://api.digitale-sammlungen.de/iiif/image/v2/bsb10538611_00286/full/1600,/0/default.jpg}{\textit{Correspondenz}, S.~278}, auf 57.4.
\par\smallskip\noindent Die Werkausgabe zeigt mit zwei Strichen an, dass sie Wisniewskys Bedeckungstafel auslässt. Die Tafel steht auf S.~307--308 der \textit{Anzeigen}. KI-Übertragung und KI-Satz.
\clearpage
